% !TEX root = ../main.tex
% ══════════════════════════════════════════════════════════════
\chapter{Choix algorithmiques et stratégies de performance}
% ══════════════════════════════════════════════════════════════
\begingroup
\let\clearpage\relax
\startcontents[chapters]
\printcontents[chapters]{}{1}{\section*{Plan}}
\endgroup
\newpage

\section*{Introduction}
\justifying
\phantomsection
\addcontentsline{toc}{section}{Introduction}

Ce chapitre présente les choix algorithmiques et les stratégies de performance qui structurent la logique métier de la plateforme. Pour chaque composant, les alternatives disponibles sont comparées sur la base de critères précis, puis la méthode retenue est justifiée.

% ─────────────────────────────────────────────────────────────
\section{Stratégie de performance base de données}
% ─────────────────────────────────────────────────────────────

Cette section présente les enjeux de performance liés à la croissance
des données, évalue les deux approches disponibles, puis justifie
la stratégie retenue.

\subsection{Contexte et enjeux}

Plusieurs tables du projet — mobilités, affectations, imports —
grossissent d'une campagne à l'autre. Le système doit pouvoir les
filtrer rapidement par année universitaire, que ce soit pour afficher
le tableau de bord RI, calculer les durées CTI ou lancer l'algorithme
d'affectation. Deux approches ont été évaluées pour répondre à cet
enjeu.

\subsection{Alternatives évaluées}

\subsubsection{Partitionnement par année}

Le partitionnement consiste à découper physiquement une table en
sous-tables, une par valeur de la colonne de partition — ici, par
année universitaire. PostgreSQL ne lit alors que la partition
concernée lors d'une requête filtrée.

\textit{Avantages :}
\begin{itemize}
  \item Réduction significative du volume scanné pour les requêtes
        filtrées par année
  \item Efficace lorsque les tables atteignent plusieurs millions
        de lignes
\end{itemize}

\textit{Limites :}
\begin{itemize}
  \item La colonne de partitionnement doit être incluse dans toutes
        les contraintes d'unicité, ce qui peut créer des conflits
        avec les migrations Django au fil des sprints
  \item Complexité supplémentaire du schéma pour un volume qui
        ne le justifie pas
\end{itemize}

\subsubsection{Indexation ciblée}

L'indexation ciblée consiste à créer des index sur les colonnes
les plus utilisées dans les filtres, sans modifier la structure
physique des tables.

\textit{Avantages :}
\begin{itemize}
  \item Aucune modification du schéma ni des migrations Django
  \item Index déclarés directement dans les modèles, visibles
        dans le code et versionnés avec lui
  \item Transparent pour l'ORM, sans configuration PostgreSQL séparée
\end{itemize}

\textit{Limites :}
\begin{itemize}
  \item Moins efficace que le partitionnement pour des volumes
        de plusieurs millions de lignes — sans incidence ici
\end{itemize}

\subsection{Synthèse comparative}

Le tableau~\ref{tab:comparaison_perf_bdd} récapitule les critères évalués pour chacune des
deux approches, afin de faciliter la comparaison et de justifier
la stratégie retenue.

\renewcommand{\arraystretch}{1.35}
\begin{table}[H]
\centering
\footnotesize
\begin{tabular}{|p{5.5cm}|>{\centering\arraybackslash}p{3cm}|>{\centering\arraybackslash}p{3cm}|}
\hline
\rowcolor{gray!20}
\textbf{Critère} & \textbf{Partitionnement} & \textbf{Indexation ciblée} \\
\hline
Adapté au volume du projet & \xmark & \cmark \\
\hline
Compatible avec les migrations Django & Contraintes d'unicité incompatibles avec certaines migrations & \cmark \\
\hline
Déclaration dans les modèles & \xmark & \cmark \\
\hline
Aucune configuration PostgreSQL séparée & \xmark & \cmark \\
\hline
Efficacité à très grande échelle & \cmark & Moins efficace pour plusieurs millions de lignes \\
\hline
\end{tabular}
\caption{Comparaison des stratégies de performance base de données}
\label{tab:comparaison_perf_bdd}
\end{table}

\subsection{Décision}

\textbf{L'indexation ciblée} est retenue. L'ENSEEIHT traite quelques
centaines d'étudiants en mobilité par an. Même en projetant ce volume
sur dix ans, les tables restent très loin du seuil où le partitionnement
apporterait un gain mesurable. Des index sont créés sur
\texttt{academic\_year\_id} et sur la combinaison
\texttt{(academic\_year\_id, student\_id)} pour les tables les plus
sollicitées, directement dans les modèles Django concernés.

% ─────────────────────────────────────────────────────────────
\section{Choix de l'algorithme d'affectation}
% ─────────────────────────────────────────────────────────────

Cette section présente les contraintes auxquelles l'algorithme doit
répondre, évalue les approches disponibles, puis justifie le choix retenu.

\subsection{Besoins spécifiques}

L'algorithme doit attribuer à chaque étudiant une destination parmi ses
vœux, en respectant un ensemble de contraintes institutionnelles. Les
résultats ont un impact réel sur les parcours : le service RI doit pouvoir
justifier chaque affectation. L'équité entre candidats et la
reproductibilité des résultats sont des exigences essentielles.

Le classement des étudiants repose principalement sur le GPA, complété
par des règles institutionnelles. Les affectations doivent respecter
les quotas par accord de mobilité et par département, les contraintes
de financement, ainsi que la situation des étudiants étrangers dans
le classement global. Chaque étudiant exprime jusqu'à trois vœux
ordonnés ; l'algorithme traite les vœux dans l'ordre et recherche
une affectation alternative en cas d'indisponibilité.

\subsection{Alternatives évaluées}

Deux approches algorithmiques ont été évaluées pour répondre à ces besoins.

\subsubsection{Algorithme glouton}

L'algorithme glouton parcourt les candidats dans l'ordre de leur GPA
et affecte chacun à son meilleur vœu encore disponible, sans jamais
revenir sur une décision déjà prise.

\textit{Avantages :}
\begin{itemize}
  \item Simple à implémenter et à expliquer
  \item Résultat prévisible, complexité $O(n^2)$
\end{itemize}

\textit{Limites :}
\begin{itemize}
  \item Le résultat dépend de l'ordre de traitement : deux étudiants
        avec des GPA très proches peuvent se retrouver dans des
        situations très différentes selon leur position dans la liste
  \item Une légère variation de l'ordre d'entrée peut produire des
        résultats différents, ce qui est incompatible avec l'exigence
        de reproductibilité
\end{itemize}

\subsubsection{Gale-Shapley}

L'algorithme de Gale-Shapley est un algorithme
d'affectation stable qui procède par rondes successives de propositions
et de rejets différés jusqu'à convergence. Il garantit qu'aucune paire
étudiant-destination ne peut former une meilleure association que celle
retenue à l'issue du processus~\cite{GaleShapley1962}.

\textit{Avantages :}
\begin{itemize}
  \item Résultat stable : il est mathématiquement impossible qu'un
        étudiant et une destination forment une meilleure paire que
        celle retenue, quelles que soient les données
  \item Équitable : le résultat ne dépend pas de l'ordre de traitement,
        deux exécutions avec les mêmes données donnent toujours le
        même résultat
  \item Optimal pour les étudiants dans la variante Student-Proposing :
        chacun obtient la meilleure destination possible compte tenu
        de son classement et de ses vœux
  \item Complexité $O(n^2)$, tout à fait adaptée à quelques centaines
        d'étudiants
\end{itemize}

\textit{Limites :}
\begin{itemize}
  \item Légèrement plus complexe à implémenter que l'algorithme glouton
\end{itemize}

\subsection{Synthèse comparative}

Le tableau~\ref{tab:comparaison_algo} récapitule les critères évalués pour chacune des deux
approches, afin de faciliter la comparaison et de justifier le choix retenu.

\renewcommand{\arraystretch}{1.35}
\begin{table}[H]
\centering
\footnotesize
\begin{tabular}{|p{5.5cm}|>{\centering\arraybackslash}p{3cm}|>{\centering\arraybackslash}p{3cm}|}
\hline
\rowcolor{gray!20}
\textbf{Critère} & \textbf{Algo. glouton} & \textbf{Gale-Shapley} \\
\hline
Stabilité des affectations & \xmark & \cmark \\
\hline
Reproductibilité & \xmark & \cmark \\
\hline
Indépendance à l'ordre de traitement & \xmark & \cmark \\
\hline
Optimalité pour les étudiants & Résultat non garanti selon l'ordre de traitement & \cmark \\
\hline
Explicabilité des résultats & Difficile à justifier si l'ordre a influencé le résultat & \cmark \\
\hline
Simplicité d'implémentation & \cmark & Plus complexe que l'algo glouton, bien documenté \\
\hline
\end{tabular}
\caption{Comparaison des algorithmes d'affectation}
\label{tab:comparaison_algo}
\end{table}

\subsection{Décision}

\textbf{Gale-Shapley} est retenu, dans sa variante \textbf{Student-Proposing} :
ce sont les étudiants qui font des propositions aux destinations, et non
l'inverse. Cette variante garantit que chaque étudiant obtient la meilleure
destination possible compte tenu de son classement et de ses vœux —
propriété dite d'\textit{optimalité côté proposant}.

Les contraintes de quota sont vérifiées à deux niveaux simultanément à
chaque proposition : un quota global par accord bilatéral (capacité de
l'université partenaire) et un quota local par département pédagogique
(places allouées au sein de l'ENSEEIHT). Une destination est disponible
seulement si les deux niveaux sont respectés ; sinon la proposition est
rejetée et l'étudiant passe à son vœu suivant.

L'algorithme glouton est écarté car il est sensible à l'ordre de traitement
et ne garantit pas la reproductibilité. Gale-Shapley assure un résultat
stable, équitable et reproductible dans ce cadre multi-contraint.


% ─────────────────────────────────────────────────────────────
\section{Choix de la méthode de répartition des quotas par département}
% ─────────────────────────────────────────────────────────────

Cette section présente les besoins liés à la répartition des quotas entre
départements, évalue les méthodes disponibles, puis justifie le choix retenu.

\subsection{Besoins spécifiques}

Chaque accord bilatéral fixe un nombre total de places disponibles. Ce nombre
doit ensuite être distribué entre les départements pédagogiques, ce qui soulève
une question qui n'a pas de réponse évidente : sur quelle base décider combien
de places revient à chaque département ?

Une répartition basée uniquement sur les demandes de l'année en cours serait
insuffisante : elle ne tient pas compte du fait qu'un département peut avoir été
systématiquement défavorisé les années précédentes, ni du fait que certains
accords voient leur nombre de places évoluer d'une campagne à l'autre. Sans
intégrer cet historique, la répartition risque d'être perçue comme arbitraire
et de creuser des déséquilibres entre départements au fil des années.

Le besoin est donc de disposer d'une méthode qui distribue les places de façon
proportionnelle à la demande réelle, tout en restant simple à comprendre,
transparente et vérifiable par le service RI. La méthode retenue doit également
garantir que la somme des places attribuées aux départements est exactement égale
au nombre de places de l'accord, sans qu'aucune place ne soit perdue lors des
arrondis.

\subsection{Alternatives évaluées}

Trois méthodes de répartition proportionnelle ont été comparées. Toutes
partent du même principe : chaque département reçoit une quote-part
proportionnelle à son nombre d'étudiants. La différence porte sur la
manière de distribuer les places entières lorsque la division produit
des résultats décimaux.

\subsubsection{Méthode Hamilton (plus grands restes)}

La méthode Hamilton attribue d'abord à chaque département la partie entière
de sa quote-part. Les places restantes sont ensuite données, une par une,
aux départements ayant les fractions les plus élevées, jusqu'à épuiser le
total disponible.

\textit{Avantages :}
\begin{itemize}
  \item Simple à comprendre et à expliquer à tous les acteurs
  \item Résultat facile à vérifier manuellement
  \item Respecte les proportions au plus près du total disponible
\end{itemize}

\textit{Limites :}
\begin{itemize}
  \item Un département dont la quote-part est inférieure à 1 ne reçoit
        aucune place garantie, même si sa fraction est élevée
\end{itemize}

\subsubsection{Méthode Jefferson (plus grandes moyennes)}

La méthode Jefferson procède de façon itérative : à chaque étape, la
prochaine place est attribuée au département dont le rapport entre sa
taille et le nombre de places déjà reçues est le plus élevé.

\textit{Avantages :}
\begin{itemize}
  \item Aucun arrondi explicite nécessaire
  \item Résultats stables et cohérents d'une campagne à l'autre
\end{itemize}

\textit{Limites :}
\begin{itemize}
  \item Avantage les grands départements et pénalise les petits
  \item Calcul itératif moins intuitif à vérifier manuellement
\end{itemize}

\subsubsection{Méthode Webster}

La méthode Webster est proche de Jefferson, mais elle utilise un arrondi
standard (seuil à 0,5) pour décider si un département reçoit une place
supplémentaire ou non.

\textit{Avantages :}
\begin{itemize}
  \item Plus équilibrée que Jefferson pour les petits départements
  \item Bien adaptée aux cas où peu de places sont à distribuer
\end{itemize}

\textit{Limites :}
\begin{itemize}
  \item Calcul moins intuitif qu'Hamilton, plus difficile à vérifier
        sans outil
  \item Résultats moins prévisibles dans les cas limites
\end{itemize}

\subsection{Synthèse comparative}

Le tableau~\ref{tab:comparaison_quotas} récapitule les critères évalués pour chacune des
trois méthodes, afin de faciliter la comparaison et de justifier le choix retenu.

\renewcommand{\arraystretch}{1.35}
\begin{table}[H]
\centering
\footnotesize
\begin{tabular}{|p{5cm}|>{\centering\arraybackslash}p{2.8cm}|>{\centering\arraybackslash}p{2.8cm}|>{\centering\arraybackslash}p{2.8cm}|}
\hline
\rowcolor{gray!20}
\textbf{Critère} & \textbf{Hamilton} & \textbf{Jefferson} & \textbf{Webster} \\
\hline
Simplicité de calcul & \cmark & \xmark & \xmark \\
\hline
Vérification manuelle possible & \cmark & Calcul itératif difficile à suivre & Moins prévisible dans les cas limites \\
\hline
Équité entre départements & \cmark & Favorise les grands départements & \cmark \\
\hline
Respect du total disponible & \cmark & \cmark & \cmark \\
\hline
Adapté à un petit nombre de places & \cmark & \xmark & \cmark \\
\hline
Transparence pour les acteurs & \cmark & \xmark & \xmark \\
\hline
\end{tabular}
\caption{Comparaison des méthodes de répartition des quotas}
\label{tab:comparaison_quotas}
\end{table}

\subsection{Décision}

\textbf{La méthode Hamilton} est retenue. Elle offre le meilleur équilibre
entre équité, simplicité et transparence dans ce contexte. Les résultats
peuvent être vérifiés à la main par le service RI, ce qui est important
pour un processus qui touche directement les chances de départ des étudiants.

La méthode Jefferson est écartée car elle avantage systématiquement les
grands départements, ce qui irait à l'encontre de l'équité visée. Webster
est plus difficile à vérifier sans outil de calcul. Hamilton est donc la
méthode la plus adaptée pour une répartition claire et justifiable.

% ─────────────────────────────────────────────────────────────
\section{Choix du moteur de recommandation}
% ─────────────────────────────────────────────────────────────

Cette section présente les besoins du module de recommandation,
évalue les approches disponibles, puis justifie la méthode retenue.

\subsection{Besoins spécifiques}

Le moteur de recommandation a pour objectif d'aider les étudiants à
choisir leurs destinations avant la saisie des vœux dans MoveON. Il
ne prend pas de décision, mais propose un classement des destinations
selon leur pertinence pour le profil de l'étudiant. Les recommandations
doivent s'appuyer sur le profil académique et sur les données historiques
disponibles, tout en restant interprétables pour l'étudiant.

\subsection{Alternatives évaluées}

Trois grandes familles d'approches ont été évaluées pour répondre à
ce besoin.

\subsubsection{Approche basée sur des règles (rule-based)}

Cette approche repose sur des règles métier simples définies manuellement,
telles que filtrer les destinations selon le département et classer par
taux d'acceptation historique.

\textit{Avantages :}
\begin{itemize}
  \item Simple à mettre en place, sans données d'apprentissage
  \item Facile à comprendre et à expliquer
\end{itemize}

\textit{Limites :}
\begin{itemize}
  \item Peu flexible, ne s'adapte pas aux profils individuels
  \item Résultats souvent trop génériques
\end{itemize}

\subsubsection{Approche par similarité (collaborative filtering)}

Cette approche propose des destinations en se basant sur les choix
d'étudiants ayant des profils similaires dans les promotions précédentes.

\textit{Avantages :}
\begin{itemize}
  \item Prend en compte les expériences passées des promotions
  \item Permet des recommandations plus personnalisées
\end{itemize}

\textit{Limites :}
\begin{itemize}
  \item Nécessite un volume important de données historiques
  \item Moins efficace pour les nouvelles destinations
  \item Mise en œuvre plus complexe
\end{itemize}

\subsubsection{Approche par apprentissage supervisé}

Cette approche consiste à entraîner un modèle pour prédire une probabilité
de compatibilité entre un étudiant et une destination. Deux modèles ont
été comparés au sein de cette famille.

\paragraph{Régression logistique}

La régression logistique est un modèle statistique qui estime la
probabilité d'un résultat à partir de variables d'entrée numériques.
Elle est particulièrement adaptée aux volumes de données modérés~\cite{scikit_learn}.

\textit{Avantages :}
\begin{itemize}
  \item Modèle simple et rapide à entraîner
  \item Résultats faciles à interpréter et à expliquer
  \item Adapté à des volumes de données modérés
\end{itemize}

\textit{Limites :}
\begin{itemize}
  \item Capture difficilement des relations complexes entre variables
\end{itemize}

\paragraph{XGBoost}

XGBoost est un algorithme de gradient boosting qui construit un ensemble
d'arbres de décision de façon séquentielle pour améliorer la précision
prédictive~\cite{scikit_learn}.

\textit{Avantages :}
\begin{itemize}
  \item Capable de modéliser des relations complexes entre variables
  \item Très bonnes performances prédictives
\end{itemize}

\textit{Limites :}
\begin{itemize}
  \item Plus complexe à entraîner et à maintenir
  \item Résultats moins interprétables
  \item Risque de surapprentissage avec un volume de données limité
\end{itemize}

\subsection{Synthèse comparative}

Le tableau~\ref{tab:comparaison_reco} récapitule les critères évalués pour chacune des
approches, afin de faciliter la comparaison et de justifier le choix retenu.

\renewcommand{\arraystretch}{1.35}
\begin{table}[H]
\centering
\footnotesize
\begin{tabular}{|p{3.8cm}|>{\centering\arraybackslash}p{2cm}|>{\centering\arraybackslash}p{2cm}|>{\centering\arraybackslash}p{2.5cm}|>{\centering\arraybackslash}p{2cm}|}
\hline
\rowcolor{gray!20}
\textbf{Critère} & \textbf{Règles} & \textbf{Collab.} & \textbf{Régression log.} & \textbf{XGBoost} \\
\hline
Personnalisation & \xmark & \cmark & \cmark & \cmark \\
\hline
Volume de données requis & \cmark & \xmark & Efficace à partir de 500 observations & Risque de surapprentissage avec peu de données \\
\hline
Interprétabilité & \cmark & Difficile sans connaître les profils similaires & \cmark & \xmark \\
\hline
Risque de surapprentissage & \cmark & Dépend de la qualité des données historiques & \cmark & \xmark \\
\hline
Simplicité de maintenance & \cmark & \xmark & \cmark & \xmark \\
\hline
\end{tabular}
\caption{Comparaison des approches de recommandation}
\label{tab:comparaison_reco}
\end{table}

\subsection{Décision}

\textbf{L'apprentissage supervisé par régression logistique} est retenu.
Le modèle estime, pour chaque couple étudiant–destination, une probabilité
de compatibilité à partir des variables suivantes :

\renewcommand{\arraystretch}{1.3}
\begin{table}[H]
\centering
\footnotesize
\begin{tabular}{|p{3.5cm}|p{8cm}|}
\hline
\rowcolor{gray!20}
\textbf{Variable} & \textbf{Description} \\
\hline
GPA de l'étudiant & Note académique normalisée, principal signal de classement \\
\hline
Département pédagogique & Filtrage sur les accords ouverts au département de l'étudiant \\
\hline
Destination & Identifiant de l'université partenaire \\
\hline
Historique des affectations & Taux de sélection par destination sur les trois dernières années \\
\hline
\end{tabular}
\caption{Variables d'entrée du moteur de recommandation}
\label{tab:features_reco}
\end{table}

Les destinations sont classées par ordre de probabilité décroissante
afin de proposer une liste de recommandations personnalisées.

La régression logistique est retenue car elle offre un bon compromis
entre simplicité, interprétabilité et facilité de maintenance. Le volume
de données — environ 900 observations sur trois ans — ne justifie pas
XGBoost, dont les performances supplémentaires s'accompagnent d'un risque
accru de surapprentissage.

\textbf{Risque de démarrage à froid.} Certaines destinations rares ou
récentes disposent de peu d'observations historiques. Pour ces cas, le
modèle se rabat sur le GPA et le département comme seuls signaux, ce qui
reste cohérent avec les règles de sélection. Le module sera activé à
partir de la deuxième année de fonctionnement, une fois qu'un historique
minimal d'une campagne complète aura été constitué.


% ─────────────────────────────────────────────────────────────
\section{Choix de l'orchestration RAG}
% ─────────────────────────────────────────────────────────────

Cette section présente les besoins du module RAG et le modèle de langage
imposé, évalue les frameworks d'orchestration disponibles, puis justifie
le choix retenu.

\subsection{Besoins spécifiques}

Le composant RAG du projet est un chatbot destiné aux étudiants, capable
de répondre à des questions sur les procédures de mobilité à partir de
la documentation institutionnelle de l'ENSEEIHT : guides, critères
d'éligibilité, délais et démarches administratives. Le cas d'usage est
relativement simple : une base documentaire stable, des questions en
français ou en anglais, et des réponses ancrées dans les documents officiels.

\subsection{Modèle de langage}

Le modèle de langage retenu est \textbf{Mistral}~\cite{mistral_dev}
(\texttt{mistral-small-latest}), conformément aux exigences de l'ENSEEIHT.
Il est accessible via l'\textbf{API Mistral ILAS}, une infrastructure
nationale dédiée aux établissements de l'enseignement supérieur français,
hébergée sur des serveurs localisés en France.

Ce mode d'accès institutionnel répond directement aux exigences RGPD :
les données traitées restent dans l'infrastructure nationale et ne sont
pas transmises à des tiers commerciaux. Il est néanmoins impératif que
les données personnelles des étudiants ne soient jamais incluses dans
les requêtes envoyées au modèle — les questions au chatbot sont
transmises telles que saisies, sans enrichissement avec des données
de dossier.


\subsection{Alternatives évaluées}

Quatre approches ont été évaluées pour orchestrer le pipeline RAG
autour du modèle Mistral.

\subsubsection{LangChain}

LangChain est un framework d'orchestration IA généraliste, compatible
avec un large éventail de modèles et de bases vectorielles~\cite{langchain_dev}.

\textit{Avantages :}
\begin{itemize}
  \item Grande communauté et nombreuses ressources disponibles
  \item Intégration simple avec Mistral et ChromaDB
  \item Adapté aux pipelines IA complexes et multi-étapes
\end{itemize}

\textit{Limites :}
\begin{itemize}
  \item API qui évolue fréquemment, induisant des risques de rupture
  \item Nombreuses dépendances, surcoût pour un chatbot documentaire simple
\end{itemize}

\subsubsection{LlamaIndex}

LlamaIndex est un framework spécialisé dans la recherche et la génération
à partir de documents structurés~\cite{LlamaIndex_dev}.

\textit{Avantages :}
\begin{itemize}
  \item Conçu spécifiquement pour les applications RAG documentaires
  \item Bonne intégration avec Mistral et les bases vectorielles
  \item API plus stable, prise en main plus directe
\end{itemize}

\textit{Limites :}
\begin{itemize}
  \item Moins adapté aux architectures IA multi-agents très complexes
\end{itemize}

\subsubsection{Haystack}

Haystack est un framework orienté production pour les systèmes de
recherche et de génération augmentée~\cite{haystack_dev}.

\textit{Avantages :}
\begin{itemize}
  \item Architecture modulaire adaptée aux déploiements en production
  \item Compatible avec Mistral
\end{itemize}

\textit{Limites :}
\begin{itemize}
  \item Documentation moins accessible pour des cas d'usage simples
  \item Intégration avec Mistral moins directe que LlamaIndex
\end{itemize}

\subsubsection{Approche directe (sans framework)}

L'approche directe consiste à implémenter manuellement le pipeline RAG
via des appels à l'API Mistral, sans couche d'abstraction intermédiaire.

\textit{Avantages :}
\begin{itemize}
  \item Contrôle total sur chaque étape du pipeline
  \item Peu de dépendances externes
\end{itemize}

\textit{Limites :}
\begin{itemize}
  \item Nécessite de développer manuellement toute la logique RAG
        (chunking, indexation, retrieval, construction du prompt)
  \item Maintenance plus complexe à long terme
\end{itemize}

\subsection{Synthèse comparative}

Le tableau~\ref{tab:comparaison_rag} récapitule les critères évalués pour chacune des
quatre approches, afin de faciliter la comparaison et de justifier
le choix retenu.

\renewcommand{\arraystretch}{1.35}
\begin{table}[H]
\centering
\footnotesize
\begin{tabular}{|p{4cm}|>{\centering\arraybackslash}p{2.2cm}|>{\centering\arraybackslash}p{2.2cm}|>{\centering\arraybackslash}p{2.2cm}|>{\centering\arraybackslash}p{2cm}|}
\hline
\rowcolor{gray!20}
\textbf{Critère} & \textbf{LangChain} & \textbf{LlamaIndex} & \textbf{Haystack} & \textbf{Direct} \\
\hline
Spécialisé RAG & Généraliste, non optimisé pour le RAG & \cmark & Orienté pipelines complexes & \xmark \\
\hline
Intégration Mistral & \cmark & \cmark & Possible mais moins directe & \cmark \\
\hline
Stabilité de l'API & \xmark & \cmark & \cmark & \cmark \\
\hline
Adapté à un cas simple & Surcoût de dépendances pour un cas simple & \cmark & Trop complexe pour un chatbot documentaire & Logique RAG à développer entièrement \\
\hline
Maintenance à long terme & API instable, risque de rupture à chaque mise à jour & \cmark & \cmark & \xmark \\
\hline
\end{tabular}
\caption{Comparaison des frameworks d'orchestration RAG}
\label{tab:comparaison_rag}
\end{table}

\subsection{Décision}

\textbf{LlamaIndex} est retenu comme framework d'orchestration RAG,
avec \textbf{ChromaDB} comme base vectorielle et
\textbf{Mistral} comme modèle de langage. Cette combinaison correspond
directement au besoin du projet : construire un chatbot documentaire
simple, basé sur une base de connaissances stable. LlamaIndex fournit
les fonctionnalités nécessaires sans surcoût de complexité, avec une
API stable et une intégration directe avec Mistral et ChromaDB.

% ─────────────────────────────────────────────────────────────
\section{Synthèse des choix}
% ─────────────────────────────────────────────────────────────

\renewcommand{\arraystretch}{1.4}

\begin{longtable}{|p{3cm}|p{3.5cm}|p{8.5cm}|}
\caption{Synthèse des choix algorithmiques et de performance} \\
\hline
\rowcolor{gray!20}
\textbf{Composant} & \textbf{Méthode retenue} & \textbf{Justification principale} \\
\hline
\endfirsthead

\hline
\rowcolor{gray!20}
\textbf{Composant} & \textbf{Méthode retenue} & \textbf{Justification principale} \\
\hline
\endhead

Stratégie de performance BDD & Indexation ciblée & Solution simple et suffisante pour le volume de données du projet \\
\hline

Algorithme d'affectation & Gale-Shapley & Stabilité, équité et reproductibilité des affectations \\
\hline

Répartition des quotas & Hamilton & Méthode proportionnelle simple, transparente et vérifiable manuellement \\
\hline

Moteur de recommandation & Régression logistique & Modèle simple, interprétable et adapté à un volume de données limité \\
\hline

Orchestration RAG & LlamaIndex + ChromaDB & Solution adaptée à un chatbot documentaire simple et stable \\
\hline

Modèle de langage & Mistral & Compatible avec le besoin multilingue du chatbot étudiant \\
\hline

\end{longtable}

\section*{Conclusion}
\justifying
\phantomsection
\addcontentsline{toc}{section}{Conclusion}

Ce chapitre a présenté les choix algorithmiques et les stratégies de performance qui structurent la logique métier de la plateforme. Chaque méthode a été sélectionnée en fonction des besoins concrets du projet, en privilégiant l'équité, la reproductibilité et la simplicité d'exploitation. Ces choix définissent le comportement intelligent du système sur lequel repose le développement présenté dans les chapitres suivants.
